\subsection{\label{sec.sign.intro}Cancellations in alternating sums}

We begin with a simple binomial identity:

\begin{proposition}
[Negative hockey-stick identity]\label{prop.binom.nhs}Let $n\in\mathbb{C}$ and
$m\in\mathbb{N}$. Then,%
\begin{equation}
\sum_{k=0}^{m}\left(  -1\right)  ^{k}\dbinom{n}{k}=\left(  -1\right)
^{m}\dbinom{n-1}{m}. \label{eq.prop.binom.nhs.eq}%
\end{equation}

\end{proposition}

There are easy proofs of this proposition by induction on $m$ or by the
telescope principle (see, e.g., \cite[Exercise 4]{18f-hw2s} and \cite[Exercise
2.1.1 \textbf{(a)} and \S 7.27]{19fco}). However, let us try to prove the
proposition combinatorially. For a bijective proof, the $\left(  -1\right)
^{k}$ and $\left(  -1\right)  ^{m}$ factors would be gamestoppers, as there is
no way to get a negative number (let alone a number of variable sign) by
counting something. However, if we think of the $\left(  -1\right)  ^{k}$ as
an opportunity for cancelling addends, then we can use combinatorics pretty well:

\begin{proof}
[Combinatorial proof of Proposition \ref{prop.binom.nhs} (sketched).]We need
to prove the equality (\ref{eq.prop.binom.nhs.eq}). Both sides of this
equality are polynomial functions in $n$. Thus, we can WLOG assume that $n$ is
a positive integer (because of the polynomial identity trick that we saw in
Subsection \ref{subsec.gf.defs.cvi}). Assume this.

Set $\left[  n\right]  :=\left\{  1,2,\ldots,n\right\}  $. We shall now
introduce some notations tailored for this particular proof.

Define an \emph{acceptable set} to be a subset of $\left[  n\right]  $ that
has size $\leq m$. Clearly,%
\begin{equation}
\left(  \text{\# of acceptable sets}\right)  =\sum_{k=0}^{m}\dbinom{n}{k}
\label{pf.prop.binom.nhs.1}%
\end{equation}
(since the \# of $k$-element subsets of $\left[  n\right]  $ equals
$\dbinom{n}{k}$ for each $k\in\mathbb{Z}$). Incidentally, we note that there
is no closed form for this sum -- another instance of the phenomenon in which
alternating sums are simpler than non-alternating ones.

Define the \emph{sign} of a finite set $I$ to be $\left(  -1\right)
^{\left\vert I\right\vert }$. Then,%
\begin{align}
&  \left(  \text{the sum of the signs of all acceptable sets}\right)
\nonumber\\
&  =\sum_{k=0}^{m}\left(  -1\right)  ^{k}\dbinom{n}{k}.
\label{pf.prop.binom.nhs.2}%
\end{align}
(This can be shown just like (\ref{pf.prop.binom.nhs.1}).)

However, the sum of the signs of all acceptable sets is a sum of $1$s and
$-1$s. Let us try to cancel as many of these $1$s and $-1$s against each other
as we can, hoping that what remains will be precisely $\left(  -1\right)
^{m}\dbinom{n-1}{m}$.

In order to cancel two addends, we need to pair up two finite sets $I$ that
have opposite signs. How do we find such pairs? One way to do so is to pick
some set $I$ that does not contain the element $1$, and pair it up with
$I\cup\left\{  1\right\}  $. Alternatively, we can pick some set $I$ that
contains the element $1$, and pair it up with $I\setminus\left\{  1\right\}
$. In other words, we pair up a finite set $I$ with either $I\setminus\left\{
1\right\}  $ or $I\cup\left\{  1\right\}  $, depending on whether $1\in I$ or
$1\notin I$.

Let us do this systematically for all finite sets: For each finite set $I$, we
define the \emph{partner} of $I$ to be the set%
\[
I^{\prime}:=%
\begin{cases}
I\setminus\left\{  1\right\}  , & \text{if }1\in I;\\
I\cup\left\{  1\right\}  , & \text{if }1\notin I
\end{cases}
\ \ \ \ =I\bigtriangleup\left\{  1\right\}  ,
\]
where the notation $X\bigtriangleup Y$ means the symmetric difference $\left(
X\cup Y\right)  \setminus\left(  X\cap Y\right)  $ of two sets $X$ and $Y$ (as
in Subsection \ref{subsec.gf.defs.commrings}). It is easy to see that each
finite set $I$ satisfies $I^{\prime\prime}=I$ and $\left\vert I^{\prime
}\right\vert =\left\vert I\right\vert \pm1$, so that $\left(  -1\right)
^{\left\vert I^{\prime}\right\vert }=-\left(  -1\right)  ^{\left\vert
I\right\vert }$. Thus, if both $I$ and $I^{\prime}$ are acceptable sets, then
their contributions to the sum of the signs of all acceptable sets cancel each
other out.

This does not mean that all addends in this sum cancel. Indeed, while each
finite set has a partner, it may happen that an acceptable set has a
non-acceptable partner, and then the contribution of the former to the sum
does not get cancelled (since the partner does not contribute to the sum).
Thus, in order to see what remains of the sum after the cancellations, we need
to study the acceptable sets that have non-acceptable partners.

Fortunately, this is easy: In order for an acceptable set $I$ to have a
non-acceptable partner, it needs to satisfy $1\notin I$ and $\left\vert
I\right\vert =m$. Better yet, this is an \textquotedblleft if and only
if\textquotedblright\ statement:

\begin{statement}
\textit{Claim 1:} Let $I$ be an acceptable set. Then, the partner $I^{\prime}$
of $I$ is non-acceptable if and only if $\left(  1\notin I\text{ and
}\left\vert I\right\vert =m\right)  $.
\end{statement}

[\textit{Proof of Claim 1:} The \textquotedblleft if\textquotedblright%
\ direction is easy: If $1\notin I$ and $\left\vert I\right\vert =m$, then the
partner $I^{\prime}$ of $I$ is defined by $I^{\prime}=I\cup\left\{  1\right\}
$ and thus has size $\left\vert I^{\prime}\right\vert =\left\vert I\right\vert
+1>\left\vert I\right\vert =m$, which shows that it is non-acceptable.

It remains to prove the converse, i.e., the \textquotedblleft only
if\textquotedblright\ direction. Thus, we assume that the partner $I^{\prime}$
of $I$ is non-acceptable. We must show that $1\notin I$ and $\left\vert
I\right\vert =m$.

Since $I$ is acceptable, we have $I\subseteq\left[  n\right]  $ and
$\left\vert I\right\vert \leq m$. If we had $1\in I$, then we would have
$I^{\prime}=I\setminus\left\{  1\right\}  \subseteq I\subseteq\left[
n\right]  $ and furthermore $\left\vert I^{\prime}\right\vert \leq\left\vert
I\right\vert $ (since $I^{\prime}\subseteq I$), so that $\left\vert I^{\prime
}\right\vert \leq\left\vert I\right\vert \leq m$. This would entail that
$I^{\prime}$ is acceptable (since $I^{\prime}\subseteq\left[  n\right]  $ and
$\left\vert I^{\prime}\right\vert \leq m$), which would contradict our
assumption that $I^{\prime}$ be non-acceptable. Hence, we cannot have $1\in
I$. Thus, $1\notin I$ is proved. Hence, $I^{\prime}=I\cup\left\{  1\right\}
$. However, $1\in\left[  n\right]  $ (since $n\geq1$), and $I^{\prime}%
=I\cup\left\{  1\right\}  \subseteq\left[  n\right]  $ (since $I\subseteq
\left[  n\right]  $ and $1\in\left[  n\right]  $). Moreover, from $I^{\prime
}=I\cup\left\{  1\right\}  $, we obtain $\left\vert I^{\prime}\right\vert
=\left\vert I\cup\left\{  1\right\}  \right\vert =\left\vert I\right\vert +1$.
Now, if we had $\left\vert I\right\vert \leq m-1$, then we would obtain
$\left\vert I^{\prime}\right\vert =\left\vert I\right\vert +1\leq m$ (since
$\left\vert I\right\vert \leq m-1$), which would entail that $I^{\prime}$ is
acceptable (since $I^{\prime}\subseteq\left[  n\right]  $), which again would
contradict our assumption. Thus, we cannot have $\left\vert I\right\vert \leq
m-1$. Hence, $\left\vert I\right\vert >m-1$, so that $\left\vert I\right\vert
\geq m$ and therefore $\left\vert I\right\vert =m$ (since $\left\vert
I\right\vert \leq m$). Thus, we have shown that $1\notin I$ and $\left\vert
I\right\vert =m$. This proves the \textquotedblleft only if\textquotedblright%
\ direction and thus completes the proof of Claim 1.]

Now, recall our line of reasoning: We start with the sum of the signs of all
acceptable sets, and we cancel any two addends that correspond to an
acceptable set and its acceptable partner. What remains are the addends
corresponding to the acceptable sets that have non-acceptable partners.
According to Claim 1, these are precisely the acceptable sets $I$ that satisfy
$\left(  1\notin I\text{ and }\left\vert I\right\vert =m\right)  $. In other
words, these are precisely the $m$-element subsets of $\left[  n\right]  $
that do not contain $1$. In other words, these are precisely the $m$-element
subsets of $\left[  n\right]  \setminus\left\{  1\right\}  $ (since a subset
of $\left[  n\right]  $ that does not contain $1$ is the same as a subset of
$\left[  n\right]  \setminus\left\{  1\right\}  $). Thus, there are precisely
$\dbinom{n-1}{m}$ of these subsets (since $\left[  n\right]  \setminus\left\{
1\right\}  $ is an $\left(  n-1\right)  $-element set), and each of them has
sign $\left(  -1\right)  ^{m}$. Hence, there are precisely $\dbinom{n-1}{m}$
addends left in the sum after our cancellations, and each of these addends is
$\left(  -1\right)  ^{m}$. Hence,
\[
\left(  \text{the sum of the signs of all acceptable sets}\right)  =\left(
-1\right)  ^{m}\dbinom{n-1}{m}.
\]
Comparing this with (\ref{pf.prop.binom.nhs.2}), we obtain
\[
\sum_{k=0}^{m}\left(  -1\right)  ^{k}\dbinom{n}{k}=\left(  -1\right)
^{m}\dbinom{n-1}{m}.
\]
This proves Proposition \ref{prop.binom.nhs}.
\end{proof}

Let me outline how to formalize this argument without using vague notions like
\textquotedblleft cancelling\textquotedblright\ and \textquotedblleft pairing
up\textquotedblright. We let%
\[
\mathcal{A}:=\left\{  \text{acceptable sets}\right\}
\]
and%
\begin{align*}
\mathcal{X}:=  &  \left\{  \text{acceptable sets whose partner is
acceptable}\right\} \\
=  &  \left\{  I\subseteq\left[  n\right]  \ \mid\ \left\vert I\right\vert
\leq m\text{ but not }\left(  \left\vert I\right\vert =m\text{ and }1\notin
I\right)  \right\}
\end{align*}
(by Claim 1 in the above proof). Now, we define a map%
\begin{align*}
f:\mathcal{X}  &  \rightarrow\mathcal{X},\\
I  &  \mapsto I^{\prime}.
\end{align*}
This map $f$ is a bijection, since each $I\in\mathcal{X}$ satisfies
$I^{\prime\prime}=I$ and thus $I^{\prime}\in\mathcal{X}$. This bijection $f$
is furthermore \emph{sign-reversing}, meaning that $\left(  -1\right)
^{\left\vert f\left(  I\right)  \right\vert }=-\left(  -1\right)  ^{\left\vert
I\right\vert }$ for all $I\in\mathcal{X}$. We claim that this automatically
guarantees%
\begin{align*}
&  \left(  \text{the sum of the signs of all acceptable sets}\right) \\
&  =\left(  \text{the sum of the signs of all acceptable sets \textbf{not}
belonging to }\mathcal{X}\right)  .
\end{align*}
The reason for this equality is that in the sum of the signs of all acceptable
sets, the contributions of the sets that belong to $\mathcal{X}$ (that is, of
the acceptable sets that have acceptable partners) cancel each other out. This
principle is worth generalizing and stating as a lemma:

\begin{lemma}
[Cancellation principle, take 1]\label{lem.sign.cancel1}Let $\mathcal{A}$ be a
finite set. Let $\mathcal{X}$ be a subset of $\mathcal{A}$.

For each $I\in\mathcal{A}$, let $\operatorname*{sign}I$ be a real number (not
necessarily $1$ or $-1$). Let $f:\mathcal{X}\rightarrow\mathcal{X}$ be a
bijection with the property that%
\begin{equation}
\operatorname*{sign}\left(  f\left(  I\right)  \right)  =-\operatorname*{sign}%
I\ \ \ \ \ \ \ \ \ \ \text{for all }I\in\mathcal{X}.
\label{eq.lem.sign.cancel1.sr}%
\end{equation}
(Such a bijection $f$ is called \emph{sign-reversing}.) Then,%
\[
\sum_{I\in\mathcal{A}}\operatorname*{sign}I=\sum_{I\in\mathcal{A}%
\setminus\mathcal{X}}\operatorname*{sign}I.
\]

\end{lemma}

Note that we did \textbf{not} require that $f\circ f=\operatorname*{id}$ in
Lemma \ref{lem.sign.cancel1}; we only required that $f$ is a bijection. That
said, most examples that I know do have $f\circ f=\operatorname*{id}$.

\begin{proof}
[Proof of Lemma \ref{lem.sign.cancel1}.]Intuitively, this is clear: The
contributions of all $I\in\mathcal{X}$ to the sum $\sum_{I\in\mathcal{A}%
}\operatorname*{sign}I$ cancel out, to the extent they are not already zero.
However, rather than formalize this cancellation, let us give an even slicker argument:

We have%
\begin{align*}
\sum_{I\in\mathcal{X}}\operatorname*{sign}I  &  =\sum_{I\in\mathcal{X}%
}\underbrace{\operatorname*{sign}\left(  f\left(  I\right)  \right)
}_{\substack{=-\operatorname*{sign}I\\\text{(by (\ref{eq.lem.sign.cancel1.sr}%
))}}}\\
&  \ \ \ \ \ \ \ \ \ \ \ \ \ \ \ \ \ \ \ \ \left(
\begin{array}
[c]{c}%
\text{here, we have substituted }f\left(  I\right)  \text{ for }I\\
\text{in the sum, since }f:\mathcal{X}\rightarrow\mathcal{X}\text{ is a
bijection}%
\end{array}
\right) \\
&  =\sum_{I\in\mathcal{X}}\left(  -\operatorname*{sign}I\right)  =-\sum
_{I\in\mathcal{X}}\operatorname*{sign}I.
\end{align*}
Adding $\sum_{I\in\mathcal{X}}\operatorname*{sign}I$ to both sides of this
equality, we obtain $2\cdot\sum_{I\in\mathcal{X}}\operatorname*{sign}I=0$.
Hence, $\sum_{I\in\mathcal{X}}\operatorname*{sign}I=0$ (since any real number
$a$ satisfying $2a=0$ must satisfy $a=0$).

Now, $\mathcal{X}\subseteq\mathcal{A}$; hence, we can split the sum
$\sum_{I\in\mathcal{A}}\operatorname*{sign}I$ as follows:%
\[
\sum_{I\in\mathcal{A}}\operatorname*{sign}I=\underbrace{\sum_{I\in\mathcal{X}%
}\operatorname*{sign}I}_{=0}+\sum_{I\in\mathcal{A}\setminus\mathcal{X}%
}\operatorname*{sign}I=\sum_{I\in\mathcal{A}\setminus\mathcal{X}%
}\operatorname*{sign}I.
\]
This proves Lemma \ref{lem.sign.cancel1}.
\end{proof}

In the proof of Proposition \ref{prop.binom.nhs}, we applied Lemma
\ref{lem.sign.cancel1} to%
\begin{align*}
\mathcal{A}  &  =\left\{  \text{acceptable sets}\right\}
\ \ \ \ \ \ \ \ \ \ \text{and}\\
\mathcal{X}  &  =\left\{  \text{acceptable sets whose partner is
acceptable}\right\}  \ \ \ \ \ \ \ \ \ \ \text{and}\\
\operatorname*{sign}I  &  =\left(  -1\right)  ^{\left\vert I\right\vert }.
\end{align*}
But there are many other situations in which Lemma \ref{lem.sign.cancel1} can
be applied. For example, $\mathcal{A}$ can be some set of permutations, and
$\operatorname*{sign}\sigma$ can be the sign of $\sigma$ (as in Definition
\ref{def.perm.sign}).

Let us observe that Lemma \ref{lem.sign.cancel1} can be generalized. Indeed,
in Lemma \ref{lem.sign.cancel1}, we can replace \textquotedblleft real
number\textquotedblright\ by \textquotedblleft element of any $\mathbb{Q}%
$-vector space\textquotedblright\ or even by \textquotedblleft element of any
additive abelian group with the property that $2a=0$ implies $a=0$%
\textquotedblright. We cannot, however, remove this requirement entirely.
Indeed, if all the signs $\operatorname*{sign}I$ were the element
$\overline{1}$ of $\mathbb{Z}/2$, then the sign-reversing condition
(\ref{eq.lem.sign.cancel1.sr}) would hold automatically (since $\overline
{1}=-\overline{1}$ in $\mathbb{Z}/2$), but the claim of Lemma
\ref{lem.sign.cancel1} would not necessarily be true.

However, if we replace the word \textquotedblleft bijection\textquotedblright%
\ by \textquotedblleft involution with no fixed points\textquotedblright, then
Lemma \ref{lem.sign.cancel1} holds even without any requirements on the group:

\begin{lemma}
[Cancellation principle, take 2]\label{lem.sign.cancel2}Let $\mathcal{A}$ be a
finite set. Let $\mathcal{X}$ be a subset of $\mathcal{A}$.

For each $I\in\mathcal{A}$, let $\operatorname*{sign}I$ be an element of some
additive abelian group. Let $f:\mathcal{X}\rightarrow\mathcal{X}$ be an
involution (i.e., a map satisfying $f\circ f=\operatorname*{id}$) that has no
fixed points. Assume that%
\[
\operatorname*{sign}\left(  f\left(  I\right)  \right)  =-\operatorname*{sign}%
I\ \ \ \ \ \ \ \ \ \ \text{for all }I\in\mathcal{X}.
\]
Then,%
\[
\sum_{I\in\mathcal{A}}\operatorname*{sign}I=\sum_{I\in\mathcal{A}%
\setminus\mathcal{X}}\operatorname*{sign}I.
\]

\end{lemma}

\begin{proof}
The idea is that all addends corresponding to the $I\in\mathcal{X}$ cancel out
from the sum $\sum_{I\in\mathcal{A}}\operatorname*{sign}I$ (because they come
in pairs of addends with opposite signs). See Section
\ref{sec.details.sign.intro} for a detailed proof.
\end{proof}

A more general version of Lemma \ref{lem.sign.cancel2} allows for $f$ to have
fixed points, as long as these fixed points have sign $0$:

\begin{lemma}
[Cancellation principle, take 3]\label{lem.sign.cancel3}Let $\mathcal{A}$ be a
finite set. Let $\mathcal{X}$ be a subset of $\mathcal{A}$.

For each $I\in\mathcal{A}$, let $\operatorname*{sign}I$ be an element of some
additive abelian group. Let $f:\mathcal{X}\rightarrow\mathcal{X}$ be an
involution (i.e., a map satisfying $f\circ f=\operatorname*{id}$). Assume that%
\[
\operatorname*{sign}\left(  f\left(  I\right)  \right)  =-\operatorname*{sign}%
I\ \ \ \ \ \ \ \ \ \ \text{for all }I\in\mathcal{X}.
\]
Assume furthermore that%
\[
\operatorname*{sign}I=0\ \ \ \ \ \ \ \ \ \ \text{for all }I\in\mathcal{X}%
\text{ satisfying }f\left(  I\right)  =I.
\]
Then,%
\[
\sum_{I\in\mathcal{A}}\operatorname*{sign}I=\sum_{I\in\mathcal{A}%
\setminus\mathcal{X}}\operatorname*{sign}I.
\]

\end{lemma}

\begin{proof}
This is similar to Lemma \ref{lem.sign.cancel2}, except that the addends
corresponding to the $I\in\mathcal{X}$ satisfying $f\left(  I\right)  =I$
don't cancel (but are already zero and thus can be removed right away). See
Section \ref{sec.details.sign.intro} for a detailed proof.
\end{proof}

Let us try to use this idea in another setting. Recall the notion of
$q$-binomial coefficients, and specifically their values (Definition
\ref{def.pars.qbinom.qbinom} \textbf{(b)}).

\begin{exercise}
\label{exe.sign.-1inom}Let $n,k\in\mathbb{N}$. Simplify $\dbinom{n}{k}_{-1}$.
\end{exercise}

\begin{example}
Let us compute $\dbinom{4}{2}_{-1}$. Theorem \ref{thm.pars.qbinom.quot1}
\textbf{(b)} yields%
\[
\dbinom{4}{2}_{q}=\dfrac{\left(  1-q^{4}\right)  \left(  1-q^{3}\right)
}{\left(  1-q^{2}\right)  \left(  1-q^{1}\right)  }.
\]
We cannot substitute $-1$ for $q$ in this formula directly, since both
numerator and denominator would become $0$ if we did. However, we can first
simplify the fraction and then substitute $-1$ for $q$: We have%
\[
\dbinom{4}{2}_{q}=\dfrac{\left(  1-q^{4}\right)  \left(  1-q^{3}\right)
}{\left(  1-q^{2}\right)  \left(  1-q^{1}\right)  }=q^{4}+q^{3}+2q^{2}+q+1,
\]
so that (by substituting $-1$ for $q$) we obtain%
\[
\dbinom{4}{2}_{-1}=\left(  -1\right)  ^{4}+\left(  -1\right)  ^{3}+2\left(
-1\right)  ^{2}+\left(  -1\right)  +1=2.
\]

\end{example}

\begin{proof}
[Solution of Exercise \ref{exe.sign.-1inom} (sketched).]Proposition
\ref{prop.pars.qbinom.alt-defs} \textbf{(b)} yields%
\[
\dbinom{n}{k}_{q}=\sum_{\substack{S\subseteq\left\{  1,2,\ldots,n\right\}
;\\\left\vert S\right\vert =k}}q^{\operatorname*{sum}S-\left(  1+2+\cdots
+k\right)  },
\]
where $\operatorname*{sum}S$ denotes the sum of the elements of a finite set
$S$ of integers. Substituting $-1$ for $q$ in this equality, we find%
\[
\dbinom{n}{k}_{-1}=\sum_{\substack{S\subseteq\left\{  1,2,\ldots,n\right\}
;\\\left\vert S\right\vert =k}}\left(  -1\right)  ^{\operatorname*{sum}%
S-\left(  1+2+\cdots+k\right)  }.
\]
Using the shorthand $\left[  n\right]  $ for the set $\left\{  1,2,\ldots
,n\right\}  $, we can rewrite this as%
\begin{equation}
\dbinom{n}{k}_{-1}=\sum_{\substack{S\subseteq\left[  n\right]  ;\\\left\vert
S\right\vert =k}}\left(  -1\right)  ^{\operatorname*{sum}S-\left(
1+2+\cdots+k\right)  }. \label{sol.sign.-1inom.1}%
\end{equation}
Let us analyze the sum on the right hand side using sign-reversing
involutions. Thus, we set%
\[
\mathcal{A}:=\left\{  S\subseteq\left[  n\right]  \ \mid\ \left\vert
S\right\vert =k\right\}  =\left\{  k\text{-element subsets of }\left[
n\right]  \right\}
\]
and%
\[
\operatorname*{sign}S:=\left(  -1\right)  ^{\operatorname*{sum}S-\left(
1+2+\cdots+k\right)  }\ \ \ \ \ \ \ \ \ \ \text{for every }S\in\mathcal{A}.
\]
Hence, (\ref{sol.sign.-1inom.1}) rewrites as%
\begin{equation}
\dbinom{n}{k}_{-1}=\sum_{S\in\mathcal{A}}\operatorname*{sign}S.
\label{sol.sign.-1inom.2}%
\end{equation}
Now, we seek a reasonable subset $\mathcal{X}\subseteq\mathcal{A}$ and a
sign-reversing bijection $f:\mathcal{X}\rightarrow\mathcal{X}$ in order to
cancel addends in the sum $\sum_{S\in\mathcal{A}}\operatorname*{sign}S$.

To wit, let us try to construct $f$ as a partial map first, and then (as an
afterthought) define $\mathcal{X}$ to be the set of all $S\in\mathcal{A}$ for
which $f\left(  S\right)  $ is defined.

Consider a $k$-element subset $S$ of $\left[  n\right]  $. What is a way to
transform $S$ that leaves its size $\left\vert S\right\vert =k$ unchanged, but
flips its sign (i.e., flips the parity of $\operatorname*{sum}S$) ? One thing
we can do is \emph{switching }$1$\emph{ with }$2$. By this I mean the
following operation:

\begin{itemize}
\item If $1\in S$ and $2\notin S$, then we replace $1$ by $2$ in $S$.

\item Otherwise, if $2\in S$ and $1\notin S$, then we replace $2$ by $1$ in
$S$.

\item If none of $1$ and $2$ is in $S$, or if both are in $S$, then we leave
$S$ unchanged for now.
\end{itemize}

\noindent Thus, switching $1$ with $2$ means replacing $S$ by%
\[
\operatorname*{switch}\nolimits_{1,2}\left(  S\right)  :=%
\begin{cases}
\left(  S\setminus\left\{  1\right\}  \right)  \cup\left\{  2\right\}  , &
\text{if }1\in S\text{ and }2\notin S;\\
\left(  S\setminus\left\{  2\right\}  \right)  \cup\left\{  1\right\}  , &
\text{if }1\notin S\text{ and }2\in S;\\
S, & \text{otherwise.}%
\end{cases}
\]
For example,%
\begin{align*}
\operatorname*{switch}\nolimits_{1,2}\left(  \left\{  1,3,5\right\}  \right)
&  =\left\{  2,3,5\right\}  ;\\
\operatorname*{switch}\nolimits_{1,2}\left(  \left\{  2,3,5\right\}  \right)
&  =\left\{  1,3,5\right\}  ;\\
\operatorname*{switch}\nolimits_{1,2}\left(  \left\{  1,2,5\right\}  \right)
&  =\left\{  1,2,5\right\}  ;\\
\operatorname*{switch}\nolimits_{1,2}\left(  \left\{  3,4,5\right\}  \right)
&  =\left\{  3,4,5\right\}  .
\end{align*}


Notice that the definition of switching $1$ with $2$ can be restated in a
somewhat slicker way using symmetric differences (see Subsection
\ref{subsec.gf.defs.commrings} for the definition of symmetric differences):%
\[
\operatorname*{switch}\nolimits_{1,2}\left(  S\right)  :=%
\begin{cases}
S\bigtriangleup\left\{  1,2\right\}  , & \text{if }\left\vert S\cap\left\{
1,2\right\}  \right\vert =1;\\
S, & \text{otherwise.}%
\end{cases}
\]
Indeed, the condition \textquotedblleft$\left\vert S\cap\left\{  1,2\right\}
\right\vert =1$\textquotedblright\ is equivalent to having either $\left(
1\in S\text{ and }2\notin S\right)  $ or $\left(  1\notin S\text{ and }2\in
S\right)  $; and in this case, the symmetric difference $S\bigtriangleup
\left\{  1,2\right\}  $ is precisely the set we need (i.e., the set $\left(
S\setminus\left\{  1\right\}  \right)  \cup\left\{  2\right\}  $ if we have
$\left(  1\in S\text{ and }2\notin S\right)  $, and the set $\left(
S\setminus\left\{  2\right\}  \right)  \cup\left\{  1\right\}  $ if we have
$\left(  1\notin S\text{ and }2\in S\right)  $).

This map $\operatorname*{switch}\nolimits_{1,2}:\mathcal{A}\rightarrow
\mathcal{A}$ is certainly a bijection (and, in fact, an involution). It is not
sign-reversing on the entire set $\mathcal{A}$; however, it has the property
that the sign of $\operatorname*{switch}\nolimits_{1,2}\left(  S\right)  $ is
opposite to the sign of $S$ whenever we have $\left(  1\in S\text{ and
}2\notin S\right)  $ or $\left(  1\notin S\text{ and }2\in S\right)  $
(because in these two cases, $\operatorname*{sum}S$ either increases by $1$ or
decreases by $1$, respectively). We can restate this property as follows: The
sign of $\operatorname*{switch}\nolimits_{1,2}\left(  S\right)  $ is opposite
to the sign of $S$ whenever we have $\left\vert S\cap\left\{  1,2\right\}
\right\vert =1$. Thus, we can use $\operatorname*{switch}\nolimits_{1,2}$ to
cancel many addends from our sum $\sum_{S\in\mathcal{A}}\operatorname*{sign}%
S$. Still, many other addends (of different signs) remain, and the result is
far from simple.

Thus, we need a \textquotedblleft Plan B\textquotedblright\ if the map
$\operatorname*{switch}\nolimits_{1,2}$ does not succeed. Assuming that
$\left\vert S\cap\left\{  1,2\right\}  \right\vert \neq1$ (that is, the set
$S\in\mathcal{A}$ contains none or both of $1$ and $2$), we gain nothing by
switching $1$ with $2$ in $S$, but maybe we get lucky switching $2$ with $3$
in $S$ (which is defined in the same way as switching $1$ with $2$, but with
the obvious changes)? If that, too, fails, we can try to switch $3$ with $4$.
If that fails as well, we can try to switch $4$ with $5$, and so on, until we
get to the end of the set $\left[  n\right]  $.

In other words, we try to define a bijection $f:\mathcal{A}\rightarrow
\mathcal{A}$ as follows: For any $S\in\mathcal{A}$, we pick the
\textbf{smallest} $i\in\left[  n-1\right]  $ such that $\left\vert
S\cap\left\{  i,i+1\right\}  \right\vert =1$ (in other words, the
\textbf{smallest} $i\in\left[  n-1\right]  $ such that exactly one of the two
elements $i$ and $i+1$ belongs to $S$); and we switch $i$ with $i+1$ in $S$
(that is, we replace $S$ by $S\bigtriangleup\left\{  i,i+1\right\}  $). This
produces a new subset $S^{\prime}$ of $\left[  n\right]  $ that has the same
size as $S$ but has the opposite sign (actually, we have $\operatorname*{sum}%
S^{\prime}=\operatorname*{sum}S\pm1$), except for the two cases when
$S=\varnothing$ and when $S=\left[  n\right]  $ (these are the cases where we
cannot find any $i\in\left[  n-1\right]  $ such that $\left\vert S\cap\left\{
i,i+1\right\}  \right\vert =1$). We set $f\left(  S\right)  :=S\bigtriangleup
\left\{  i,i+1\right\}  $.

Here are some examples (for $n=4$ and $k=2$):%
\begin{align*}
f\left(  \left\{  1,3\right\}  \right)   &  =\left\{  2,3\right\}
\ \ \ \ \ \ \ \ \ \ \left(  \text{here, the smallest }i\text{ is }1\right)
;\\
f\left(  \left\{  1,4\right\}  \right)   &  =\left\{  2,4\right\}
\ \ \ \ \ \ \ \ \ \ \left(  \text{here, the smallest }i\text{ is }1\right)
;\\
f\left(  \left\{  3,4\right\}  \right)   &  =\left\{  2,4\right\}
\ \ \ \ \ \ \ \ \ \ \left(  \text{here, the smallest }i\text{ is }2\right)  .
\end{align*}
Alas, the last two of these examples show that $f$ is not injective (as
$f\left(  \left\{  1,4\right\}  \right)  =\left\{  2,4\right\}  =f\left(
\left\{  3,4\right\}  \right)  $). Thus, $f$ is not a bijection. The
underlying problem is that the $i$ that was picked in the construction of
$f\left(  S\right)  $ is not uniquely recoverable from $f\left(  S\right)  $.
Hence, our map $f$ does not work for us -- we cannot use it to cancel addends,
since we cannot cancel (e.g.) a single $1$ against multiple $-1$s.

How can we salvage this argument? We change our map $f$ to \textquotedblleft
space\textquotedblright\ our switches apart. That is, we again start by trying
to switch $1$ with $2$; if this fails, we jump straight to trying to switch
$3$ with $4$; if this fails too, we jump further to trying to switch $5$ with
$6$; and so on, until we either succeed at some switch or run out of pairs to
switch. For the explicit description of $f$, this means that instead of
picking the \textbf{smallest} $i\in\left[  n-1\right]  $ such that $\left\vert
S\cap\left\{  i,i+1\right\}  \right\vert =1$, we pick the \textbf{smallest
odd} $i\in\left[  n-1\right]  $ such that $\left\vert S\cap\left\{
i,i+1\right\}  \right\vert =1$; and then we set $f\left(  S\right)
:=S\bigtriangleup\left\{  i,i+1\right\}  $ as before.

In other words, we define our new map $f:\mathcal{A}\rightarrow\mathcal{A}$ as
follows: For any $S\in\mathcal{A}$, we set%
\[
f\left(  S\right)  :=S\bigtriangleup\left\{  i,i+1\right\}  ,
\]
where $i$ is the \textbf{smallest odd} element of $\left[  n-1\right]  $ such
that $\left\vert S\cap\left\{  i,i+1\right\}  \right\vert =1$. If no such $i$
exists, we just set $f\left(  S\right)  :=S$. (We will soon see when this happens.)

Here are some examples (for $n=8$ and $k=3$):%
\begin{align*}
f\left(  \left\{  1,3,4\right\}  \right)   &  =\left\{  2,3,4\right\}
\ \ \ \ \ \ \ \ \ \ \left(  \text{here, the smallest odd }i\text{ is
}1\right)  ;\\
f\left(  \left\{  2,4,5\right\}  \right)   &  =\left\{  1,4,5\right\}
\ \ \ \ \ \ \ \ \ \ \left(  \text{here, the smallest odd }i\text{ is
}1\right)  ;\\
f\left(  \left\{  1,2,3\right\}  \right)   &  =\left\{  1,2,4\right\}
\ \ \ \ \ \ \ \ \ \ \left(  \text{here, the smallest odd }i\text{ is
}3\right)  ;\\
f\left(  \left\{  3,5,7\right\}  \right)   &  =\left\{  4,5,7\right\}
\ \ \ \ \ \ \ \ \ \ \left(  \text{here, the smallest odd }i\text{ is
}3\right)  ;\\
f\left(  \left\{  3,4,5\right\}  \right)   &  =\left\{  3,4,6\right\}
\ \ \ \ \ \ \ \ \ \ \left(  \text{here, the smallest odd }i\text{ is
}5\right)  ;\\
f\left(  \left\{  5,6,7\right\}  \right)   &  =\left\{  5,6,8\right\}
\ \ \ \ \ \ \ \ \ \ \left(  \text{here, the smallest odd }i\text{ is
}7\right)  .
\end{align*}
And here are two more examples (for $n=8$ and $k=4$):%
\begin{align*}
f\left(  \left\{  1,2,5,7\right\}  \right)   &  =\left\{  1,2,6,7\right\}
\ \ \ \ \ \ \ \ \ \ \left(  \text{here, the smallest odd }i\text{ is
}5\right)  ;\\
f\left(  \left\{  1,2,5,6\right\}  \right)   &  =\left\{  1,2,5,6\right\}
\ \ \ \ \ \ \ \ \ \ \left(  \text{here, there is no appropriate odd }i\right)
.
\end{align*}


Once again, it is clear that the set $f\left(  S\right)  $ has size $k$
whenever $S$ does. Hence, $f:\mathcal{A}\rightarrow\mathcal{A}$ is at least a
well-defined map. This time, the map $f$ is furthermore an involution (that
is, $f\circ f=\operatorname*{id}$). Here is a quick argument for this (details
are left to the reader): Since we have \textquotedblleft
spaced\textquotedblright\ the switches apart, they don't interfere with each
other. Thus, the $i$ that gets chosen in the construction of $f\left(
S\right)  $ will again get chosen in the construction of $f\left(  f\left(
S\right)  \right)  $ (since the elements of $f\left(  S\right)  $ that are
smaller than this $i$ will not have changed from $S$). Thus, the switch that
happens in the construction of $f\left(  f\left(  S\right)  \right)  $ undoes
the switch made in the construction of $f\left(  S\right)  $, and as a result,
the set $f\left(  f\left(  S\right)  \right)  $ will be $S$ again. This shows
that $f\circ f=\operatorname*{id}$.

Thus, $f$ is an involution, hence a bijection. Moreover, $\operatorname*{sign}%
\left(  f\left(  S\right)  \right)  =-\operatorname*{sign}S$ holds whenever
$f\left(  S\right)  \neq S$ (because $f\left(  S\right)  \neq S$ implies that
$f\left(  S\right)  =S\bigtriangleup\left\{  i,i+1\right\}  $ for some $i$
satisfying $\left\vert S\cap\left\{  i,i+1\right\}  \right\vert =1$, and
therefore $\operatorname*{sum}\left(  f\left(  S\right)  \right)
=\operatorname*{sum}S\pm1$). Thus, we set%
\[
\mathcal{X}:=\left\{  S\in\mathcal{A}\ \mid\ f\left(  S\right)  \neq
S\right\}  ,
\]
and we restrict $f$ to a map $\mathcal{X}\rightarrow\mathcal{X}$ (this is
well-defined, since it is easy to see from $f\circ f=\operatorname*{id}$ that
$f\left(  S\right)  \in\mathcal{X}$ for each $S\in\mathcal{X}$). Then, the map
$f$ becomes a sign-reversing bijection from $\mathcal{X}$ to $\mathcal{X}$.
Hence, Lemma \ref{lem.sign.cancel1} yields%
\[
\sum_{I\in\mathcal{A}}\operatorname*{sign}I=\sum_{I\in\mathcal{A}%
\setminus\mathcal{X}}\operatorname*{sign}I.
\]
Renaming the index $I$ as $S$, we can rewrite this equality as%
\[
\sum_{S\in\mathcal{A}}\operatorname*{sign}S=\sum_{S\in\mathcal{A}%
\setminus\mathcal{X}}\operatorname*{sign}S.
\]
Hence, (\ref{sol.sign.-1inom.2}) becomes%
\begin{equation}
\dbinom{n}{k}_{-1}=\sum_{S\in\mathcal{A}}\operatorname*{sign}S=\sum
_{S\in\mathcal{A}\setminus\mathcal{X}}\operatorname*{sign}S.
\label{sol.sign.-1inom.3}%
\end{equation}


Now, what is $\mathcal{A}\setminus\mathcal{X}$ ? In other words, what addends
are left behind uncancelled?

In order to answer this question, we need to consider the case when $n$ is
even and the case when $n$ is odd separately. We begin with the case when $n$
is even.

A $k$-element subset $S$ of $\left[  n\right]  $ belongs to $\mathcal{A}%
\setminus\mathcal{X}$ if and only if it satisfies $f\left(  S\right)  =S$. In
other words, $S$ belongs to $\mathcal{A}\setminus\mathcal{X}$ if and only if
there exists no odd $i\in\left[  n-1\right]  $ such that $\left\vert
S\cap\left\{  i,i+1\right\}  \right\vert =1$ (because $f$ has been defined in
such a way that $f\left(  S\right)  =S$ in this case, while $f\left(
S\right)  =S\bigtriangleup\left\{  i,i+1\right\}  \neq S$ in the other case).
In other words, $S$ belongs to $\mathcal{A}\setminus\mathcal{X}$ if and only
if for each odd $i\in\left[  n-1\right]  $, the size $\left\vert S\cap\left\{
i,i+1\right\}  \right\vert $ is either $0$ or $2$. This is equivalent to
saying that if we break up the $n$ elements $1,2,\ldots,n$ into $n/2$
\textquotedblleft blocks\textquotedblright\
\[
\left\{  1,2\right\}  ,\ \ \left\{  3,4\right\}  ,\ \ \left\{  5,6\right\}
,\ \ \ldots,\ \ \left\{  n-1,n\right\}
\]
(this can be done, since $n$ is even), then the intersection of $S$ with each
block has size $0$ or $2$. In other words, this is saying that the set $S$
consists of entire blocks (i.e., each block is either fully included in $S$ or
is disjoint from $S$). In other words, this is saying that the set $S$ is a
union (possibly empty) of blocks. We call a subset $S$ of $\left[  n\right]  $
\emph{blocky} if it satisfies this condition.\footnote{For example, $\left\{
3,4,5,6,9,10\right\}  $ is a blocky subset of $\left[  10\right]  $, whereas
$\left\{  2,3,5,6\right\}  $ is not (since it neither fully includes nor is
disjoint from the block $\left\{  1,2\right\}  $).} Thus, a $k$-element subset
$S$ of $\left[  n\right]  $ belongs to $\mathcal{A}\setminus\mathcal{X}$ if
and only if it is blocky. In other words, $\mathcal{A}\setminus\mathcal{X}$ is
the set of all blocky $k$-element subsets of $\left[  n\right]  $.

How many blocky $k$-element subsets does $\left[  n\right]  $ have, and what
are their signs? Any blocky subset of $\left[  n\right]  $ has the
form\footnote{The symbol \textquotedblleft$\sqcup$\textquotedblright\ means
\textquotedblleft disjoint union\textquotedblright\ (in our case, a union of
disjoint sets).}%
\[
\left\{  i_{1},i_{1}+1\right\}  \sqcup\left\{  i_{2},i_{2}+1\right\}
\sqcup\cdots\sqcup\left\{  i_{p},i_{p}+1\right\}
\]
for some odd elements $i_{1}<i_{2}<\cdots<i_{p}$ of $\left[  n-1\right]  $.
Thus, this subset has size $2p$, which entails that its size is even. Hence,
if $k$ is odd, there are no blocky $k$-element subsets of $\left[  n\right]  $
at all. In other words, if $k$ is odd, then there are no $S\in\mathcal{A}%
\setminus\mathcal{X}$ (since $\mathcal{A}\setminus\mathcal{X}$ is the set of
all blocky $k$-element subsets of $\left[  n\right]  $). Therefore, if $k$ is
odd, then the sum $\sum_{S\in\mathcal{A}\setminus\mathcal{X}}%
\operatorname*{sign}S$ is empty, and thus (\ref{sol.sign.-1inom.3}) simplifies
to%
\begin{equation}
\dbinom{n}{k}_{-1}=\sum_{S\in\mathcal{A}\setminus\mathcal{X}}%
\operatorname*{sign}S=\left(  \text{empty sum}\right)  =0.
\label{sol.sign.-1inom.ev/od}%
\end{equation}


Let us now consider the case when $k$ is even. In this case, again, any blocky
subset $S$ of $\left[  n\right]  $ has the form%
\[
\left\{  i_{1},i_{1}+1\right\}  \sqcup\left\{  i_{2},i_{2}+1\right\}
\sqcup\cdots\sqcup\left\{  i_{p},i_{p}+1\right\}
\]
for some odd elements $i_{1}<i_{2}<\cdots<i_{p}$ of $\left[  n-1\right]  $.
Moreover, again, this subset has size $2p$. Thus, if $S$ has size $k$, then we
must have $2p=k$, so that $p=k/2$. Thus, any blocky $k$-element subset $S$ of
$\left[  n\right]  $ has the form
\[
\left\{  i_{1},i_{1}+1\right\}  \sqcup\left\{  i_{2},i_{2}+1\right\}
\sqcup\cdots\sqcup\left\{  i_{k/2},i_{k/2}+1\right\}
\]
for some odd elements $i_{1}<i_{2}<\cdots<i_{k/2}$ of $\left[  n-1\right]  $.
Hence, there are $\dbinom{n/2}{k/2}$ such subsets $S$ (since there are
precisely $\dbinom{n/2}{k/2}$ choices for these odd elements $i_{1}%
<i_{2}<\cdots<i_{k/2}$\ \ \ \ \footnote{because the set $\left[  n-1\right]  $
has $n/2$ odd elements}). Moreover, any such subset $S$ satisfies%
\begin{align*}
\operatorname*{sum}S  &  =\underbrace{i_{1}+\left(  i_{1}+1\right)  }%
_{\equiv1\operatorname{mod}2}+\underbrace{i_{2}+\left(  i_{2}+1\right)
}_{\equiv1\operatorname{mod}2}+\cdots+\underbrace{i_{k/2}+\left(
i_{k/2}+1\right)  }_{\equiv1\operatorname{mod}2}\\
&  \equiv\underbrace{1+1+\cdots+1}_{k/2\text{ times}}=k/2\operatorname{mod}2
\end{align*}
and therefore%
\[
\underbrace{\operatorname*{sum}S}_{\equiv k/2\operatorname{mod}2}%
-\underbrace{\left(  1+2+\cdots+k\right)  }_{=\dfrac{k\left(  k+1\right)  }%
{2}}\equiv k/2-\dfrac{k\left(  k+1\right)  }{2}=-k^{2}/2\equiv
0\operatorname{mod}2
\]
(since $k$ is even, so that $-k^{2}/2$ is even), and%
\begin{equation}
\operatorname*{sign}S=\left(  -1\right)  ^{\operatorname*{sum}S-\left(
1+2+\cdots+k\right)  }=1 \label{sol.sign.-1inom.od/ev.sign}%
\end{equation}
(since $\operatorname*{sum}S-\left(  1+2+\cdots+k\right)  \equiv
0\operatorname{mod}2$). Hence, the sum $\sum_{S\in\mathcal{A}\setminus
\mathcal{X}}\operatorname*{sign}S$ has $\dbinom{n/2}{k/2}$ addends (because
$\mathcal{A}\setminus\mathcal{X}$ is the set of all blocky $k$-element subsets
of $\left[  n\right]  $, and we have just shown that there are $\dbinom
{n/2}{k/2}$ such subsets), and each of these addends is $1$ (by
(\ref{sol.sign.-1inom.od/ev.sign})). Thus, this sum simplifies to%
\[
\sum_{S\in\mathcal{A}\setminus\mathcal{X}}\operatorname*{sign}S=\dbinom
{n/2}{k/2}\cdot1=\dbinom{n/2}{k/2}.
\]
Hence, (\ref{sol.sign.-1inom.3}) becomes%
\begin{equation}
\dbinom{n}{k}_{-1}=\sum_{S\in\mathcal{A}\setminus\mathcal{X}}%
\operatorname*{sign}S=\dbinom{n/2}{k/2}. \label{sol.sign.-1inom.ev/ev}%
\end{equation}


Thus, we have computed $\dbinom{n}{k}_{-1}$

\begin{itemize}
\item in the case when $n$ is even and $k$ is odd (obtaining
(\ref{sol.sign.-1inom.ev/od})), and

\item in the case when $n$ is even and $k$ is even (obtaining
(\ref{sol.sign.-1inom.ev/ev})).
\end{itemize}

It remains to handle the case when $n$ is odd. This case is different in that
the $n$ elements $1,2,\ldots,n$ are now subdivided into $\left(  n+1\right)
/2$ \textquotedblleft blocks\textquotedblright%
\[
\left\{  1,2\right\}  ,\ \ \left\{  3,4\right\}  ,\ \ \left\{  5,6\right\}
,\ \ \ldots,\ \ \left\{  n-2,n-1\right\}  ,\ \ \left\{  n\right\}  ,
\]
with the last of these blocks having size $1$. As a consequence, this time, a
blocky subset of $\left[  n\right]  $ can have odd size. Moreover, the parity
of $k$ determines whether a blocky $k$-element subset of $\left[  n\right]  $
will contain $n$:

\begin{itemize}
\item If $k$ is even, then no blocky $k$-element subset of $\left[  n\right]
$ can contain $n$ (because if it did, then it would have odd size, since all
non-$\left\{  n\right\}  $ blocks have even size).

\item If $k$ is odd, then every blocky $k$-element subset of $\left[
n\right]  $ must contain $n$ (because if it didn't, then it would have even
size, since all non-$\left\{  n\right\}  $ blocks have even size).
\end{itemize}

Thus, when classifying the blocky $k$-element subsets of $\left[  n\right]  $,
we can either dismiss $n$ immediately (if $k$ is even) or take $n$ for granted
(if $k$ is odd); in either case, the problem gets reduced to classifying the
blocky $k$-element or $\left(  k-1\right)  $-element subsets of $\left[
n-1\right]  $, which we already know how to do (since $n-1$ is even). The
result is that the \# of blocky $k$-element subsets of $\left[  n\right]  $
(in the case when $n$ is odd) is%
\[%
\begin{cases}
\dbinom{\left(  n-1\right)  /2}{k/2}, & \text{if }k\text{ is even;}\\
\dbinom{\left(  n-1\right)  /2}{\left(  k-1\right)  /2}, & \text{if }k\text{
is odd}%
\end{cases}
\ \ \ \ \ =\dbinom{\left(  n-1\right)  /2}{\left\lfloor k/2\right\rfloor },
\]
and their signs are always $1$. Hence, we obtain%
\[
\sum_{S\in\mathcal{A}\setminus\mathcal{X}}\operatorname*{sign}S=\dbinom
{\left(  n-1\right)  /2}{\left\lfloor k/2\right\rfloor }\cdot1=\dbinom{\left(
n-1\right)  /2}{\left\lfloor k/2\right\rfloor }.
\]
Thus, (\ref{sol.sign.-1inom.3}) becomes%
\begin{equation}
\dbinom{n}{k}_{-1}=\sum_{S\in\mathcal{A}\setminus\mathcal{X}}%
\operatorname*{sign}S=\dbinom{\left(  n-1\right)  /2}{\left\lfloor
k/2\right\rfloor }. \label{sol.sign.-1inom.od/*}%
\end{equation}
This is the answer to our exercise in the case when $n$ is odd.

Returning to the general case, we can now combine the formulas
(\ref{sol.sign.-1inom.ev/ev}), (\ref{sol.sign.-1inom.ev/od}) and
(\ref{sol.sign.-1inom.od/*}) into a single equality that holds in all cases:%
\begin{equation}
\dbinom{n}{k}_{-1}=%
\begin{cases}
0, & \text{if }n\text{ is even and }k\text{ is odd;}\\
\dbinom{\left\lfloor n/2\right\rfloor }{\left\lfloor k/2\right\rfloor }, &
\text{otherwise.}%
\end{cases}
\label{sol.sign.-1inom.res}%
\end{equation}

\end{proof}

This formula (\ref{sol.sign.-1inom.res}) can be generalized. Indeed, here is a
generalization of the number $-1$:

\begin{definition}
\label{def.root-of-unity.prim}Let $K$ be a field. Let $d$ be a positive
integer. \medskip

\textbf{(a)} A\emph{ }$d$\emph{-th root of unity} in $K$ means an element
$\omega$ of $K$ satisfying $\omega^{d}=1$. In other words, a $d$-th root of
unity in $K$ means an element of $K$ whose $d$-th power is $1$. \medskip

\textbf{(b)} A \emph{primitive }$d$\emph{-th root of unity} in $K$ means an
element $\omega$ of $K$ satisfying
\[
\omega^{d}=1
\]
but
\[
\omega^{i}\neq1\ \ \ \ \ \ \ \ \ \ \text{for each }i\in\left\{  1,2,\ldots
,d-1\right\}  .
\]


In other words, a primitive $d$-th root of unity in $K$ means an element of
the multiplicative group $K^{\times}$ whose order is $d$.
\end{definition}

For $K=\mathbb{C}$, the $d$-th roots of unity are the $d$ complex numbers
\newline$e^{2\pi i0/d},e^{2\pi i1/d},e^{2\pi i2/d},\ldots,e^{2\pi i\left(
d-1\right)  /d}$ (which are the vertices of a regular $d$-gon inscribed in the
unit circle, with one vertex at $1$), whereas the primitive $d$-th roots of
unity are the numbers $e^{2\pi ig/d}$ for all $g\in\left[  d\right]  $
satisfying $\gcd\left(  g,d\right)  =1$. In particular, the $2$-nd roots of
unity in $\mathbb{C}$ are $1$ and $-1$, and the only primitive $2$-nd root of
unity is $-1$. The following picture shows the six $6$-th roots of unity in
$\mathbb{C}$ (with the two primitive $6$-th roots colored blue, and the
remaining four roots colored red):%
\[%
% [tikz figure removed]
%BeginExpansion
% [tikz figure removed]%
%EndExpansion
\ \ \ .
\]


We can now generalize (\ref{sol.sign.-1inom.res}) by replacing $-1$ by
primitive roots of unity:

\begin{theorem}
[$q$-Lucas theorem]\label{thm.sign.q-lucas}Let $K$ be a field. Let $d$ be a
positive integer. Let $\omega$ be a primitive $d$-th root of unity in $K$. Let
$n,k\in\mathbb{N}$. Let $q$ and $u$ be the quotients obtained when dividing
$n$ and $k$ by $d$ with remainder, and let $r$ and $v$ be the respective
remainders. Then,%
\begin{equation}
\dbinom{n}{k}_{\omega}=\dbinom{q}{u}\cdot\dbinom{r}{v}_{\omega}.
\label{eq.thm.sign.q-lucas.eq}%
\end{equation}

\end{theorem}

Note that the equality (\ref{eq.thm.sign.q-lucas.eq}) contains two $\omega
$-binomial coefficients and one regular binomial coefficient.

It is not hard to check that (\ref{sol.sign.-1inom.res}) is the particular
case of Theorem \ref{thm.sign.q-lucas} for $d=2$ and $\omega=-1$. Indeed, the
only possible remainders of an integer upon division by $2$ are $0$ and $1$,
and the $\omega$-binomial coefficients $\dbinom{r}{v}_{\omega}$ for
$r,v\in\left\{  0,1\right\}  $ are%
\[
\dbinom{0}{0}_{\omega}=1,\ \ \ \ \ \ \ \ \ \ \dbinom{0}{1}_{\omega
}=0,\ \ \ \ \ \ \ \ \ \ \dbinom{1}{0}_{\omega}=1,\ \ \ \ \ \ \ \ \ \ \dbinom
{1}{1}_{\omega}=1.
\]


Theorem \ref{thm.sign.q-lucas} can be proved using a generalization of
sign-reversing involutions to $d$-cycles instead of pairs\footnote{In our
proofs so far, we have been using the equality $1+\left(  -1\right)  =0$ to
cancel addends in alternating sums. This equality can be generalized as
follows: If $\omega$ is a primitive $d$-th root of unity for $d>1$, then
$1+\omega+\omega^{2}+\cdots+\omega^{d-1}=0$. This equality can be used to
cancel addends in sums that involve powers of $\omega$. The \emph{discrete
Fourier transform} (see, e.g., \cite[\S 5.6]{OlvSha}) and the
\emph{roots-of-unity filter} (see, e.g., \cite[\S 1.2.9.D]{Knuth-TAoCP1}) are
applications of this idea.}. A more algebraic proof -- using \textquotedblleft
noncommutative generating functions\textquotedblright\ -- is given in Exercise
\ref{exe.sign.qbinom.qlucas}.

\begin{noncompile}
TODO: Do the combinatorial proof of Theorem \ref{thm.sign.q-lucas} using a
Fourier generalization of sign-reversing involutions.
\end{noncompile}
